\section{Zariski pairs}
\label{section-zariski-pairs}

\noindent
In this section and the next a {\it pair} is a pair $(A, I)$ where $A$
is a ring and $I \subset A$ is an ideal. A {\it morphism of pairs}
$(A, I) \to (B, J)$ is a ring map $\varphi : A \to B$ with
$\varphi(I) \subset J$.

\begin{definition}
\label{definition-zariski-pair}
A {\it Zariski pair} is a pair $(A, I)$ such that
$I$ is contained in the Jacobson radical of $A$.
\end{definition}

\begin{lemma}
\label{lemma-idempotents-determined-modulo-radical}
Let $(A, I)$ be a Zariski pair. Then the map from
idempotents of $A$ to idempotents of $A/I$ is injective.
\end{lemma}

\begin{proof}
An idempotent of a local ring is either $0$ or $1$.
Thus an idempotent is determined by the set of maximal ideals
where it vanishes, by
Algebra, Lemma \ref{algebra-lemma-characterize-zero-local}.
\end{proof}

\begin{lemma}
\label{lemma-check-isomorphism-zariski}
Let $(A, I)$ be a Zariski pair. Let $A \to B$ be a flat,
integral, finitely presented ring map such that $A/I \to B/IB$
is an isomorphism. Then $A \to B$ is an isomorphism.
\end{lemma}

\begin{proof}
The ring map $A \to B$ is finite by
Algebra, Lemma \ref{algebra-lemma-characterize-finite-in-terms-of-integral}.
Hence $B$ is finitely presented as an $A$-module by
Algebra, Lemma \ref{algebra-lemma-finite-finitely-presented-extension}.
Hence $B$ is a finite locally free $A$-module by
Algebra, Lemma \ref{algebra-lemma-finite-projective}.
Since the module $B$ has rank $1$
along $V(I)$ (see rank function described in
Algebra, Lemma \ref{algebra-lemma-finite-projective}),
and as $(A, I)$ is a Zariski pair, we conclude that the
rank is $1$ everywhere.
It follows that $A \to B$ is an isomorphism:
it is a pleasant exercise to show that
a ring map $R \to S$ such that $S$ is a locally free $R$-module
of rank $1$ is an isomorphism (hint: look at local rings).
\end{proof}

\begin{lemma}
\label{lemma-helper-finite}
Let $(A, I)$ be a Zariski pair. Let $A \to B$ be a finite ring map.
Assume
\begin{enumerate}
\item $B/IB = B_1 \times B_2$ is a product of $A/I$-algebras
\item $A/I \to B_1/IB_1$ is surjective,
\item $b \in B$ maps to $(1, 0)$ in the product.
\end{enumerate}
Then there exists a monic $f \in A[x]$ with $f(b) = 0$ and
$f \bmod I = (x - 1)x^d$ for some $d \geq 1$.
\end{lemma}

\begin{proof}
By Lemma \ref{lemma-lift-idempotent-upstairs} we can find an
\'etale ring map $A \to A'$ inducing an isomorphism $A/I \to A'/IA'$
such that $B' = B \otimes_A A'$ contains
an idempotent $e'$ lifting the image of $b$ in $B'/IB'$.
Consider the corresponding $A'$-algebra decomposition
$$
B' = B'_1 \times B'_2
$$
which is compatible with the one given in the lemma upon
reduction modulo $I$.
The map $A' \to B'_1$ is surjective modulo $IA'$. By Nakayama's lemma
(Algebra, Lemma \ref{algebra-lemma-NAK})
we can find $i \in IA'$ such that after replacing $A'$
by $A'_{1 + i}$ the map $A' \to B'_1$ is surjective.
Observe that the image $b'_1 \in B'_1$ of $b$ 
satisfies $b'_1 - 1 \in IB'_1$.
Thus we may pick $a' \in IA'$ mapping to $b'_1 - 1$.
On the other hand, the image $b'_2 \in B'_2$ of $b$ is in $IB'_2$. By
Algebra, Lemma \ref{algebra-lemma-integral-integral-over-ideal}
there exist a monic polynomial
$g(x) = x^d + \sum a'_j x^j$ of degree $d$ with $a'_j \in IA'$ such that
$g(b'_2) = 0$ in $B'_2$. Thus the image $b' = (b'_1, b'_2) \in B'$
of $b$ is a root of the polynomial $(x - 1 - a')g(x)$. We conclude that
$$
(b' - 1)(b')^d \in \sum\nolimits_{j = 0, \ldots, d} IA' \cdot (b')^j
$$
We claim that this implies
$$
(b - 1)b^d \in \sum\nolimits_{j = 0, \ldots, d} I \cdot b^j
$$
in $B$. For this it is enough to see that the ring map
$A \to A'$ is faithfully flat, because the condition is that
the image of $(b - 1)b^d$ is zero in $B/\sum_{j = 0, \ldots, d} Ib^j$
(use Algebra, Lemma \ref{algebra-lemma-faithfully-flat-universally-injective}).
The map $A \to A'$ flat because it is \'etale
(Algebra, Lemma \ref{algebra-lemma-etale}).
On the other hand, the induced map on spectra is open
(see Algebra, Proposition \ref{algebra-proposition-fppf-open} and use
previous lemma referenced) and the image
contains $V(I)$. Since $I$ is contained in the Jacobson radical of $A$
we conclude.
\end{proof}

\begin{lemma}
\label{lemma-noetherian-zariski-jacobson-complement}
Let $(A, I)$ be a Zariski pair with $A$ Noetherian. Let $f \in I$.
Then $A_f$ is a Jacobson ring.
\end{lemma}

\begin{proof}
We will use the criterion of
Algebra, Lemma \ref{algebra-lemma-noetherian-dim-1-Jacobson}.
Let $\mathfrak p \subset A$ be a prime ideal such that
$\mathfrak p_f = \mathfrak p A_f$ is prime and not maximal. We have to
show that $A_f/\mathfrak p_f = (A/\mathfrak p)_f$
has infinitely many prime ideals. After replacing $A$ by $A/\mathfrak p$
we may assume $A$ is a domain, $\dim A_f > 0$, and our goal is to
show that $\Spec(A_f)$ is infinite. Since $\dim A_f > 0$ we can
find a nonzero prime ideal $\mathfrak q \subset A$ not containing $f$.
Choose a maximal ideal $\mathfrak m \subset A$ containing $\mathfrak q$.
Since $(A, I)$ is a Zariski pair, we see $I \subset \mathfrak m$.
Hence $\mathfrak m \not = \mathfrak q$ and $\dim(A_\mathfrak m) > 1$.
Hence $\Spec((A_\mathfrak m)_f) \subset \Spec(A_f)$ is infinite by
Algebra, Lemma \ref{algebra-lemma-Noetherian-local-domain-dim-2-infinite-opens}
and we win.
\end{proof}











